\section{Система поддержки принятия решений}
\subsection{Постановка задачи и формат диалога}
Система рекомендует игры на основе предпочтений и двух уточнений.
Ввод имеет фиксированную форму «Мне нравятся: RPG, инди».
Возраст не включён: исходная база не содержит рейтингов, и приложение
не может обосновать возрастную фильтрацию. Пример возраста в задании
является иллюстрацией возможных слотов, а не единственно допустимым форматом.

Словарь нормализации объединяет русские и английские названия в пять
канонических предпочтений: ролевые игры, инди, головоломки, шутеры и
пошаговые стратегии. Например, RPG, РПГ и «ролевая игра» дают один атом
\code{role_playing_game}. Словарь полностью приведён в README приложения.
Инди является производственным признаком.
Для него используется предикат indie/1; жанры проверяются отношением has\_genre/2.

Первый уточняющий вопрос: режим «один», «вместе» или «любой».
Второй: нужна ли игра для новичка («да» или «нет»).
Пустой и неизвестный ответ приводит к повтору соответствующего вопроса.
Рекомендации выводятся только после завершения диалога.

\subsection{Архитектура и алгоритм}
Python отвечает за диалог, проверку формата и отображение.
Модуль bridge.pl загружает games.pl первой лабораторной и выполняет логический
поиск. Входные параметры передаются JSON через стандартный ввод отдельного
процесса SWI-Prolog, результаты возвращаются через стандартный вывод.
Командная оболочка не используется; пользовательский текст не подставляется в код цели.

\begin{lstlisting}
INPUT preference line
PARSE complete fixed format
NORMALIZE known aliases; REMOVE duplicates
ASK mode; ASK beginner filter
FOR each game in the knowledge base:
    REQUIRE selected mode
    REQUIRE beginner rule if requested
    COLLECT matching preferences
    KEEP game if at least one preference matched
SORT by match count descending, then game identifier
OUTPUT recommendations and matched conditions
\end{lstlisting}

Предпочтения объединяются по ИЛИ, а уточнения применяются по И.
Оценка кандидата $g$ равна числу совпавших различных предпочтений:
\[
s(g)=\left|\{t\in T:\operatorname{matches}(g,t)\}\right|.
\]
Требуется $s(g)>0$. Одинаковые оценки упорядочиваются по идентификатору,
что делает результат воспроизводимым. В текущих данных комбинации жанров
и инди не пересекаются, поэтому фактические оценки в сценариях равны единице.
Счёт поддерживает расширение базы, но не изображает вероятности или качество игр.

\begin{lstlisting}
candidate(Tags, Mode, Experience, G, Score, Matched) :-
    game(G), mode_fits(G, Mode),
    experience_fits(G, Experience),
    findall(T, (member(T, Tags), matches(G, T)), Raw),
    sort(Raw, Matched), length(Matched, Score), Score > 0.
\end{lstlisting}

При $N$ играх и $K$ предпочтениях логическая фильтрация выполняет порядка $NK$
проверок совпадения; сортировка $R$ результатов требует порядка $R\log R$ сравнений.
Для 12 игр стоимость запуска интерпретатора важнее самого поиска.
Индексацию простых фактов обеспечивает SWI-Prolog; отсечение не требуется.

\subsection{Сценарии использования}
Каждая строка ниже соответствует реальному вызову моста Prolog, записанному
в evidence.json. Начало входной строки во всех случаях --- «Мне нравятся:».

\begin{longtable}{P{3.0cm}P{3.0cm}P{8.0cm}}
\caption{Входы, рекомендации и объяснения}\\\toprule
Предпочтения & Уточнения & Результат и основание\\\midrule\endfirsthead
\toprule Предпочтения & Уточнения & Результат и основание\\\midrule\endhead
RPG & один; новичок & cyberpunk\_2077, the\_witcher\_3: жанр RPG, single\_player и recommended\_for\_beginner.\\
Инди & вместе; новичок & stardew\_valley: indie, multiplayer и правило новичка.\\
Головоломка & вместе; любой опыт & portal\_2: has\_genre(puzzle) и multiplayer.\\
Шутер & один; любой опыт & doom\_eternal: has\_genre(first\_person\_shooter) и single\_player.\\
Стратегия & любой режим; любой опыт & civilization\_vi: has\_genre(turn\_based\_strategy).\\
RPG & вместе; любой опыт & Пустой список: ни у одной игры с записанным RPG нет multiplayer.\\
Шутер & один; новичок & Пустой список: Doom Eternal отмечена difficult и исключена правилом новичка.\\
RPG, инди & один; новичок & cyberpunk\_2077, stardew\_valley, the\_witcher\_3: каждый совпал с одним предпочтением.\\\bottomrule
\end{longtable}

Например, второй сценарий формирует запрос:
\begin{lstlisting}
candidate([indie],multi,beginner,G,S,M)
\end{lstlisting}
Для Stardew Valley проверяются признаки indie и multiplayer, затем правило
recommended\_for\_beginner. Итоговый балл равен 1,
а список объясняющих совпадений равен [indie].

\subsection{Ошибки и ограничения}
Неизвестное предпочтение «гонки», пустая строка и завершающая запятая отвергаются
до запуска Prolog. Повторение RPG и РПГ не увеличивает балл.
При конфликте режима и признаков система сообщает, что совпадений в БЗ нет,
и предлагает изменить режим или снять фильтр новичка.
Отсутствие swipl приводит к инструкции установить интерпретатор;
на выполнение запроса установлен таймаут 15 секунд. EOF и Ctrl+C завершают диалог.

Пустой список не доказывает отсутствие подходящей игры в мире.
В частности, жанр Baldur's Gate 3 в этой БЗ не задан, поэтому система не
рекомендует её по RPG. Такой результат согласован с данными и явно документирован.
Признак «подходит новичку» не является рекомендацией по возрасту.
